\section{Attractive determinant large deviations}
\label{att:entropy-section}

The argument here concerns the actual Slater determinant on $X=\mathbb
P^2$, with anticanonical degree $3k$, $N=\binom{3k+2}{2}$, and
$E_k=-(kN)^{-1}\log\|\Delta_k\|^2$. Write $\sigma$ for
Fubini--Study probability. The unconditioned attractive ensemble is
\[
 Z_k(\gamma)=\int e^{\gamma N E_k}\,d\sigma^{\otimes N},
 \qquad P_{k,\gamma}=Z_k(\gamma)^{-1}e^{\gamma N E_k}\sigma^{\otimes N}.
\]
The scalar inverse-section estimate ensures finiteness for $0\le\gamma<1/3$.
The proof first obtains the product-law mean energy from the established
repulsive pressure, using a uniform two-sided estimate at zero coupling.
Compact full conditional densities then give the entropy and energy
limits for the actual attractive tilted Gibbs laws. No generic
microscopic compactness hypothesis is imposed on arbitrary symmetric
laws, and no attractive pressure statement is used as an input.

\subsection{Entropy is determined by coarse types when conditionals are compact}

Let $(X,\sigma)$ be a compact probability space with an increasing
sequence $\mathcal A_j$ of finite Borel partitions generating its Borel
sigma field modulo null sets. Cells of positive probability suffice.
Let $\mathbb E_j$ denote conditional expectation onto $\mathcal A_j$.

\begin{lemma}[An exact coarse-graining bound]
\label{att:coarse-entropy}
Let $P_N\ll\sigma^{\otimes N}$ have finite relative entropy, and let
$f_{i,x_{-i}}$ be its full one-coordinate conditional density relative
to $\sigma$. Denote the law of its cell labels by $P_N^{(j)}$ and the
reference cell probabilities by $\sigma_j$. Then
\begin{align}
 0&\le D(P_N\Vert\sigma^{\otimes N})-
               D(P_N^{(j)}\Vert\sigma_j^{\otimes N})\\
 &\le\sum_{i=1}^N\mathbb E_{P_{N,-i}}
    \left[\operatorname{Ent}_\sigma(f_{i,x_{-i}})
              -\operatorname{Ent}_\sigma(\mathbb E_j f_{i,x_{-i}})\right].
 \label{att:coarse-bound}
\end{align}
Here $\operatorname{Ent}_\sigma(f)=\int f\log f\,d\sigma$ for a
probability density $f$.
\end{lemma}

\begin{proof}
Condition on the entire cell-label vector $Y$. Relative entropy chain
rule identifies the difference on the left with
\[
 \mathbb E_Y D\left(P_N(\cdot\mid Y)\,
       \middle\Vert\,\bigotimes_i\sigma(\cdot\mid Y_i)\right).
\]
Tensorize relative entropy inside this product of cells, using full
one-coordinate conditionals. Given $x_{-i}$, the conditional law of
$x_i$ within cell $A$ has density $f_i/\int_Af_i\,d\sigma$ against
$\sigma|_A$. Averaging the relative entropy to normalized $\sigma|_A$
over its cell probabilities gives exactly
\[
 \int f_i\log\frac{f_i}{\mathbb E_j f_i}\,d\sigma
 =\operatorname{Ent}_\sigma(f_i)-
                       \operatorname{Ent}_\sigma(\mathbb E_jf_i).
\]
This proves the upper bound. The lower bound is data processing.
\end{proof}

\begin{lemma}[Entropy chaos from compact conditional densities]
\label{att:entropy-chaos}
Suppose $P_N$ is exchangeable and every full one-coordinate conditional
density belongs to a fixed relatively compact subset $\mathcal F$ of
$L^r(\sigma)$ for some $r>1$. Suppose the empirical laws converge to
$\Gamma\in\mathcal P(\mathcal P(X))$, and partition boundaries have
zero mass for the limiting first marginal. Then
\[
 \lim_N\frac1N D(P_N\Vert\sigma^{\otimes N})
 =\int D(\mu\Vert\sigma)\,d\Gamma(\mu).
\]
In particular, if empirical measures converge in probability to
$\sigma$, the relative entropy per particle tends to zero.
\end{lemma}

\begin{proof}
The operators $\mathbb E_j$ converge strongly to the identity in
$L^r$ and are contractions. Convergence is uniform on the compact
closure of $\mathcal F$. The entropy functional is continuous on
$L^r$-convergent families of densities; uniform $L^r$ control supplies
uniform integrability of $f\log f$. Consequently
\[
 \epsilon_j:=\sup_{f\in\mathcal F}
 [\operatorname{Ent}_\sigma(f)-
       \operatorname{Ent}_\sigma(\mathbb E_jf)]\longrightarrow0.
\]
By Lemma~\ref{att:coarse-entropy}, the fine and coarse entropies per
particle differ by at most $\epsilon_j$.

For a fixed partition with $M$ cells, exchangeability makes the label
law uniform on every type class. If $\nu_N$ is the vector of empirical
cell frequencies, the method-of-types bounds give
\[
 \left|N^{-1}D(P_N^{(j)}\Vert\sigma_j^{\otimes N})
               -\mathbb E D(\nu_N\Vert\sigma_j)\right|
 \le \frac{2M\log(N+1)}N.
\]
One can obtain this directly from the entropy of the type law, which
is at most $M\log(N+1)$, and the corresponding upper and lower bounds
on multinomial type probabilities. The limiting frequencies have law
$(\mu\mapsto\mu_j)_*\Gamma$. Relative entropy on this finite simplex
is continuous and bounded, so the last expectation converges to
$\int D(\mu_j\Vert\sigma_j)\,d\Gamma$.
Let $N\to\infty$ and then $j\to\infty$. The finite-partition
entropies increase to full relative entropy. Tensorization also
bounds the original entropy per particle by
$\sup_{f\in\mathcal F}\operatorname{Ent}_\sigma(f)<\infty$,
so the limiting integral is finite. Monotone convergence completes
the proof.
\end{proof}

\subsection{The conditional compactness supplied by complex geometry}

\begin{lemma}[Compact attractive conditionals]
\label{att:conditional-compactness}
Fix $0\le\gamma<1/3$. The full one-coordinate conditional densities
of $P_{k,\gamma}$, for all $k$, lie in a fixed compact subset of
$L^r(\sigma)$ for some $r>1$.
\end{lemma}

\begin{proof}
For almost every retained tuple the determinant in the remaining
coordinate is a nonzero section $s$ of $\mathcal O(3k)$. Put
\[
 u=k^{-1}\log\|s\|^2-\sup_X k^{-1}\log\|s\|^2.
\]
Then $u$ is $\omega_0$-plurisubharmonic and $\sup u=0$. Its
conditional density is $e^{-\gamma u}/\int e^{-\gamma u}\,d\sigma$.
The normalized $\omega_0$-psh family is compact in $L^1$.
For every $0\le a<1/3$ it satisfies the uniform exponential estimate
\[
 \int e^{-a u}\,d\sigma\le C_a,
 \qquad C_a=\frac2{(1-3a)(2-3a)}.
\]
Choose $1<r<r'$ with $r'\gamma<1/3$. Every sequence of normalized
potentials has an $L^1$-convergent subsequence and, on a further
subsequence, pointwise almost-everywhere convergence. The exponential
bound at $r'\gamma$ and Vitali's theorem imply convergence of
$e^{-\gamma u}$ in $L^r$. Its integral is at least one, and
normalization is continuous in $L^r$. Thus the normalized densities
have compact closure in $L^r$. This argument includes limits with
analytic poles. At $\gamma=0$ the family consists of the constant one.
\end{proof}

The same conclusion holds with a fixed bounded continuous one-body
field, or with a smooth cylindrical empirical tilt whose one-coordinate
increments form a fixed compact bounded family. Such an extension
asserts compactness of conditionals; it does not assert empirical
concentration or identify a pressure for an arbitrary tilt.

\subsection{A uniform quadratic bound near zero attraction}

\begin{lemma}[A small-coupling exponential estimate]
\label{att:quadratic-pressure}
Let $\nu=g\sigma$ be a probability measure with $g\in L^p(\sigma)$
for some $p>1$. There are $t_0>0$ and $C_\nu<\infty$, independent
of $k$, such that
\begin{equation}
 t\,\mathbb E_{\nu^{\otimes N}}E_k
 \le \frac1N\log\int e^{tN E_k}\,d\nu^{\otimes N}
 \le t\,\mathbb E_{\nu^{\otimes N}}E_k+C_\nu t^2,
 \qquad |t|\le t_0.
 \label{att:quadratic-bound}
\end{equation}
The constants can be chosen uniformly when the exponent $p>1$ and the $L^p$ bound of $g$ are fixed. The assertion includes finiteness of the integrals and mean energies.
\end{lemma}

\begin{proof}
Let $p'=p/(p-1)$, and choose $s>0$ with $p's<1/3$. For every
normalized $\omega_0$-psh $u\le0$, H\"older gives
\[
 \int e^{-su}\,d\nu\le
 M:=\|g\|_p C_{p's}^{1/p'}<\infty.
\]
The section inverse estimate can be anchored at any nonzero value,
since its normalized maximum dominates that value. Applying it
successively in all particle coordinates shows that the joint
partition is finite for $t\le s$, at every level. Its logarithmic
moments are finite on compact subintervals below $s$.

For a full conditional at parameter $t$, set $Y=-u\ge0$ and
$F(t)=\log\int e^{tY}\,d\nu$. Its relative entropy to $\nu$ is
$tF'(t)-F(t)$. Jensen and the exponential bound give
$\mathbb E_\nu Y\le s^{-1}\log M$. Hence, for $|t|\le s/2$,
\[
 \int e^{tY}\,d\nu\ge e^{t\mathbb E_\nu Y}\ge M^{-1/2},
 \qquad
 F''(t)\le M^{1/2}\int Y^2e^{sY/2}\,d\nu
                         \le\frac{16M^{3/2}}{s^2}.
\]
In particular
\[
 0\le tF'(t)-F(t)=\int_0^t rF''(r)\,dr
                         \le C_\nu t^2,
 \qquad C_\nu=8M^{3/2}/s^2.
\]
For negative $t$, the integral has negative integrand and reversed
orientation, so it is still nonnegative. Tensorization of relative
entropy for the full Gibbs law yields
$D(P_{k,t}^{\nu}\Vert\nu^{\otimes N})/N\le C_\nu t^2$ for both signs.
Repulsive integrability follows also from the Hadamard lower bound;
attractive integrability on a larger interval was already shown.

Write $p_k(t)=N^{-1}\log Z_k^\nu(t)$. On either side of zero,
\[
 \frac{d}{dt}\frac{p_k(t)}t
 =\frac{t p_k'(t)-p_k(t)}{t^2}
 =\frac{D(P_{k,t}^{\nu}\Vert\nu^{\otimes N})}{Nt^2}
 \in[0,C_\nu].
\]
Integrating between zero and $t$, and multiplying by $t$, gives
\[
 0\le p_k(t)-t\mathbb E_{\nu^{\otimes N}}E_k\le C_\nu t^2
\]
for both signs, since both integration orientation and multiplication
reverse when $t<0$. Differentiability at zero follows from the two-sided
exponential neighborhood. This proves the claimed estimate with
$t_0=s/2$.
\end{proof}

\begin{theorem}[Product-law mean energy from repulsive pressure]
\label{att:general-product-energy}
For every smooth strictly positive probability density $\nu$ on
$\mathbb P^2$,
\[
 \mathbb E_{\nu^{\otimes N_k}}E_k\longrightarrow E(\nu).
\]
Together with the deterministic lower comparison, this convergence
also holds in $L^1$.
\end{theorem}

\begin{proof}
For $a>0$, the established unrestricted repulsive determinant LDP
\cite[Theorems 1.1 and 5.7]{prob_berman2017}, with volume reference
$\nu$ and the original metric, gives
\[
 -\frac1{aN}\log Z_k^\nu(-a)\longrightarrow
 J_\nu(a):=\inf_\mu\left[E(\mu)+\frac1a D(\mu\Vert\nu)\right].
\]
The energy normalization is unchanged by changing the volume reference;
only the entropy reference changes. If the repulsive theorem is stated
using a $\nu$-orthonormal basis, its Gram matrix relative to the fixed
$\sigma$-orthonormal basis has spectrum between $\min g$ and $\max g$.
The resulting additive change of $E_k$ is $O(k^{-1})$, so its limit
is unchanged. Lemma~\ref{att:quadratic-pressure}
gives, for $0<a\le t_0$,
\[
 J_\nu(a)\le\liminf_k\mathbb E_{\nu^N}E_k
 \le\limsup_k\mathbb E_{\nu^N}E_k\le J_\nu(a)+C_\nu a.
\]
We claim $J_\nu(a)\to E(\nu)$ as $a\downarrow0$. The upper bound
uses $\mu=\nu$. Approximate minimizers have relative entropy at most
$a(E(\nu)+o(1))$, since $E\ge0$, so they converge in total variation
to $\nu$. Lower semicontinuity of $E$ gives the reverse limiting
inequality. Letting $a\downarrow0$ proves convergence of the means.
For the $L^1$ statement, empirical convergence to $\nu$ and the
continuous-envelope lower comparison make negative deviations from
$E(\nu)$ tend to zero in probability. They are bounded by
$E(\nu)+\log N/k$. Convergence of the means then forces both deviations
to vanish in mean.
\end{proof}

The theorem is obtained from the repulsive sign, not from an assumed
negative-temperature limit. Its proof applies more generally whenever
the polarized manifold has the stated uniform normalized-section
exponential estimate and the repulsive determinant pressure theorem.
The separate toric assignment formula gives a direct constructive
calculation of the same product mean energy on invariant reference laws.

\subsection{Where the limiting empirical densities lie}

Fix $0<\gamma<1/3$. Let $\Phi$ be Lipschitz on $\mathcal P(X)$ for
the Wasserstein distance of any fixed smooth metric on $X$, and put
\[
 dP_{k,\gamma,\Phi}=Z_{k,\gamma,\Phi}^{-1}
       e^{N[\gamma E_k-\Phi(L_N)]}\,d\sigma^N.
\]
Finiteness and finite relative entropy follow from the strict scalar
integrability margin and boundedness of $\Phi$.

\begin{lemma}[A compact class for all directing densities]
\label{att:directing-class}
For these laws there are $r>1$, $c>0$, and a compact convex set
$\mathcal K\subset L^r(\sigma)$ of probability densities, all at least
$c$, with the following properties. Every full one-coordinate conditional
density belongs to $\mathcal K$. Every subsequential empirical limit
$\Pi$ is supported on the measures $g\sigma$ with $g\in\mathcal K$.
Weak probability topology and $L^r$ topology coincide on this compact set.
\end{lemma}

\begin{proof}
The conditional exponent consists of $-\gamma u(x)$ as in
Lemma~\ref{att:conditional-compactness}, plus
\[
 -V_{k,x_{-i}}(x),\qquad
 V_{k,x_{-i}}(x)=N[\Phi(L_N^x)-\Phi(L_N^{x_*})].
\]
If $L$ is a Lipschitz constant of $\Phi$, these functions have absolute
value at most $L\operatorname{diam}X$ and Lipschitz constant at most
$L$, independently of all parameters. Their closure is compact in
$C(X)$ by Arzel\`a--Ascoli. The psh compactness proof therefore makes
the normalized conditional densities a relatively compact family in
$L^r$, for any $r>1$ with $r\gamma<1/3$. Their normalization is at
most $e^{L\operatorname{diam}X}C_\gamma$, whereas their numerators are
at least $e^{-L\operatorname{diam}X}$. Hence the common lower bound
$c=e^{-2L\operatorname{diam}X}/C_\gamma$ holds.

Take the closed convex hull in $L^r$. It is compact: approximate the
original compact set by a finite epsilon-net; its convex hull lies in
the epsilon-neighborhood of a compact finite-dimensional polytope.
The norm bound, normalization, and positive lower bound persist.

Here is an elementary support argument, avoiding any assumption about
conditional de Finetti limits. For $h\in C(X)$ set
$b_h=\sup_{g\in\mathcal K}\int hg\,d\sigma$. A one-coordinate
conditional given previously revealed coordinates is a mixture of full
conditionals and therefore still belongs to $\mathcal K$. Thus
$\mathbb E[h(X_i)\mid X_1,\ldots,X_{i-1}]\le b_h$. Hoeffding's
bounded-increment supermartingale estimate gives
\[
 P_{k,\gamma,\Phi}(L_N h>b_h+\epsilon)
 \le \exp[-2N\epsilon^2/(\operatorname{osc}h)^2]
\]
when $h$ is nonconstant. Apply this to a countable dense collection of
continuous functions. Separation of a compact convex set of probability
measures shows that every limiting empirical measure belongs to
$\mathcal K$. Finally the continuous injection
$\mathcal K\subset L^r\to\mathcal P(X)$ is a homeomorphism onto
its image, by compactness and the Hausdorff property.
\end{proof}

\begin{lemma}[Uniform smooth reference approximation]
\label{att:smooth-net}
For every $\epsilon>0$, the compact set $\mathcal K$ above has a
finite set of smooth positive reference densities $g_1,\ldots,g_M$
such that every $g\in\mathcal K$ is assigned a reference satisfying
\[
 D(g\sigma\Vert g_a\sigma)<\epsilon,
 \qquad |E(g\sigma)-E(g_a\sigma)|<\epsilon.
\]
The references can always have the same positive lower bound $c$ and
a common $L^r$ bound, independently of $\epsilon$.
\end{lemma}

\begin{proof}
Apply the heat semigroup of a unitary-invariant metric. It preserves
probability, positivity, and the lower bound, and contracts $L^r$.
Strong continuity is uniform on compact sets. Its small-time images,
together with $\mathcal K$, have compact closure in $L^r$.
Relative entropy is jointly continuous on this closure: the map
$g\mapsto\log g$ is continuous into $L^{r'}$ since $g\ge c$, while
logarithmic growth is controlled by arbitrary small positive powers
of $g$. H\"older then treats the cross-entropy term, and
$g\mapsto\int g\log g$ is continuous by $L^r$ integrability.

For energy, normalized psh potentials have a uniform $L^{r'}$ bound
by their exponential integrability. The energy dual consequently gives
\[
 |E(g\sigma)-E(h\sigma)|\le C_{r'}\|g-h\|_r.
\]
This bound follows by taking the difference of the two suprema in the
energy dual and normalizing each test potential to have supremum zero.
For an explicit relative-entropy estimate, choose
$1<p\le\min(r,2)$. The elementary bound
$x\log x-x+1\le C_p|x-1|^p$ for $x\ge0$ gives
\[
 D(g\sigma\Vert h\sigma)
 \le C_p c^{1-p}\|g-h\|_p^p\qquad(h\ge c).
\]
A finite net in the heat-smoothed compact set now has the required
properties.
\end{proof}

\subsection{The actual tilted Gibbs mean energy}

\begin{theorem}[Entropy and energy limits of attractive Gibbs laws]
\label{att:gibbs-energy-limit}
For fixed $0<\gamma<1/3$ and Lipschitz $\Phi$, suppose a subsequence
of the empirical laws of $P_{k,\gamma,\Phi}$ converges to $\Pi$.
Then
\begin{align}
 \frac1N D(P_{k,\gamma,\Phi}\Vert\sigma^N)
       &\longrightarrow\int H(\mu)\,d\Pi(\mu),\\
 \mathbb E_{P_{k,\gamma,\Phi}}E_k
       &\longrightarrow\int E(\mu)\,d\Pi(\mu).
 \label{att:gibbs-energy-limit-eq}
\end{align}
The statements concern the actual attractive Gibbs laws. They do not
assert microscopic energy compactness for all symmetric laws of
bounded entropy and energy.
\end{theorem}

\begin{proof}
Lemma~\ref{att:directing-class} and the entropy-chaos lemma give the
entropy limit. Partition boundaries for the underlying spatial
partitions may be chosen $\sigma$-null, since the limiting first
marginal has an $L^r$ density.

Choose finite smooth references from Lemma~\ref{att:smooth-net}.
Because weak and strong topology agree on $\mathcal K$, finitely many
weak neighborhoods can be chosen so that every point of
$\mathcal K$ in a given neighborhood has the prescribed small relative
entropy and energy error to its reference. Choose radii with
$\Pi$-null boundaries and make the neighborhoods into disjoint cells
$A_1,\ldots,A_M$. Add a remaining cell $A_0$ with reference
$\nu_0=\sigma$; its limiting probability is zero. Put
$p_{k,a}=P_{k,\gamma,\Phi}(L_N\in A_a)$, and let $R_{k,a}$ be the
corresponding conditional law whenever its probability is nonzero.

The disjoint-event entropy identity and change of reference give
\begin{align*}
 \sum_a\frac{p_{k,a}}N D(R_{k,a}\Vert\nu_a^N)
 &=\frac1N D(P_{k,\gamma,\Phi}\Vert\sigma^N)
   +\frac{H((p_{k,a})_a)}N\\
 &\quad-\sum_a p_{k,a}\mathbb E_{R_{k,a}}
                    \int\log g_a\,dL_N.
\end{align*}
For this fixed finite family every $\log g_a$ is bounded continuous.
The entropy limit and convergence on the cells imply
\begin{equation}
 \limsup_k\sum_a\frac{p_{k,a}}N
                D(R_{k,a}\Vert\nu_a^N)
 \le\sum_a\int_{A_a}D(\mu\Vert\nu_a)\,d\Pi(\mu)<\epsilon.
 \label{att:mixture-entropy}
\end{equation}
There is no need for a separate estimate on the exceptional cell:
its nonnegative entropy term is included in this sum.

The two-sided pressure bound and the entropy variational inequality
imply, for either sign and a common $0<t\le t_0$,
\[
 \left|\mathbb E_{R_{k,a}}E_k-
                   \mathbb E_{\nu_a^N}E_k\right|
 \le Ct+\frac{D(R_{k,a}\Vert\nu_a^N)}{Nt}.
\]
The constants $C,t_0$ are uniform over every refinement of the net:
Lemma~\ref{att:quadratic-pressure} uses only the common $L^r$ bound.
Multiply by $p_{k,a}$ and sum. Theorem~\ref{att:general-product-energy}
handles the finitely many smooth references, and
\eqref{att:mixture-entropy} gives
\[
 \limsup_k\left|\mathbb E_{P_{k,\gamma,\Phi}}E_k
                   -\int E(\mu)\,d\Pi(\mu)\right|
 \le Ct+\epsilon/t+\epsilon.
\]
First refine the reference net, then let $t\downarrow0$. This proves
mean-energy convergence. The argument controls collective logarithmic
singularities through entropy and conditional regularity; no spectral
positive-moment bound is substituted for a logarithmic estimate.
\end{proof}

\subsection{The full Laplace principle and shrinking design constraints}

Set $F_\gamma(\mu)=H(\mu)-\gamma E(\mu)$ for finite-entropy
probability measures, and $F_\gamma=+\infty$ otherwise. This convention
never subtracts two infinite quantities: finite entropy implies finite
energy by the scalar exponential estimate. The following argument also
applies on an arbitrary positively polarized compact manifold, replacing
$1/3$ by its normalized-psh integrability threshold $\alpha$.

\begin{lemma}[Good rate and entropy-continuous energy recovery]
\label{att:good-rate}
For $0\le\gamma<\alpha$, $F_\gamma$ is lower semicontinuous and
has compact sublevels. For every $\gamma<a<\alpha$,
\begin{equation}
 F_\gamma(\mu)\ge(1-\gamma/a)H(\mu)
                          -(\gamma/a)\log C_a,
 \label{att:coercivity}
\end{equation}
where $C_a$ bounds $\int e^{-au}\,d\sigma$ over all normalized psh
potentials. If $\mu_j\to\mu$ weakly and
$H(\mu_j)\to H(\mu)<\infty$, then $E(\mu_j)\to E(\mu)$.
\end{lemma}

\begin{proof}
The case $\gamma=0$ for the first assertion is ordinary entropy lower
semicontinuity. Take $0<\gamma<\alpha$. Write $\mathcal A$ for
Aubin--Mabuchi potential energy, extended by $-\infty$ on infinite-energy
normalized psh functions. The measure energy dual gives
\begin{equation}
 F_\gamma(\mu)=
 \inf_{\substack{u\in\operatorname{PSH}(X,\omega_0)\\\sup u=0}}
 \left\{D(\mu\Vert e^{-\gamma u}\sigma)
                       -\gamma\mathcal A(u)\right\}.
 \label{att:compact-potential-dual}
\end{equation}
Here relative entropy against a finite, not necessarily normalized,
measure $\nu$ means $\int\log(d\mu/d\nu)\,d\mu$ and has the
variational formula
\[
 D(\mu\Vert\nu)=\sup_{f\in C(X)}
       \left(\int f\,d\mu-\log\int e^f\,d\nu\right).
\]
For finite $H(\mu)$, the exponential bound and the entropy inequality
make $\int|u|\,d\mu$ finite, uniformly over normalized $u$.
Thus $D(\mu\Vert e^{-\gamma u}\sigma)=H(\mu)+\gamma\int u\,d\mu$,
and \eqref{att:compact-potential-dual} follows from the energy dual.
Potentials with $\mathcal A(u)=-\infty$ contribute $+\infty$.

The equality also holds at infinite entropy. To see this and obtain a
uniform bound, choose $\gamma<a<\alpha$ and put $r=a/\gamma>1$.
For every bounded continuous $f$, H\"older gives
\[
 \int e^{f/r'}e^{-\gamma u}\,d\sigma
 \le \left(\int e^f\,d\sigma\right)^{1/r'} C_a^{1/r}.
\]
Taking the relative-entropy variational supremum therefore yields
\[
 D(\mu\Vert e^{-\gamma u}\sigma)
 \ge (1-\gamma/a)H(\mu)-(\gamma/a)\log C_a.
\]
This is valid with infinite values. Since $\mathcal A(u)\le0$ when
$\sup u=0$, every objective in \eqref{att:compact-potential-dual}
satisfies the same bound. Consequently the infimum is $+\infty$ when
$H(\mu)=+\infty$, and \eqref{att:coercivity} follows in all cases.

The normalized psh family is compact in $L^1$, and $\mathcal A$ is
upper semicontinuous on it, including its value $-\infty$; these are
standard variational pluripotential inputs \cite{BBGZ2013}.
If $u_j\to u$ in this compact family, the strict exponent margin
and Vitali's theorem give
$e^{-\gamma u_j}\to e^{-\gamma u}$ in $L^1(\sigma)$.
The displayed entropy variational formula makes $D(\mu\Vert\nu)$
jointly lower semicontinuous under weak convergence of its two measure
arguments. Hence the objective in \eqref{att:compact-potential-dual}
is jointly lower semicontinuous in $(\mu,u)$. Minimizing over the
compact potential space proves lower semicontinuity of $F_\gamma$:
choose approximate minimizing potentials along any sequence and extract
a convergent subsequence. Compactness of $\mathcal P(X)$ then makes
every free-energy sublevel compact.

Finally choose any fixed $0<\gamma<\alpha$. Under the assumptions
of the last assertion, lower semicontinuity of $F_\gamma$ gives
\[
 H(\mu)-\gamma E(\mu)
 \le H(\mu)-\gamma\limsup_j E(\mu_j).
\]
Thus $\limsup E(\mu_j)\le E(\mu)$; lower semicontinuity of $E$
gives the reverse lower limit. All energies are finite and uniformly
bounded by the entropy/exponential estimate. This proves energy recovery.
\end{proof}

\begin{theorem}[Full attractive determinant Laplace principle]
\label{att:full-laplace}
For every fixed $0\le\gamma<1/3$ and every
$\Phi\in C(\mathcal P(\mathbb P^2))$,
\begin{equation}
 \lim_{k\to\infty}\frac1N\log\int
       e^{N[\gamma E_k-\Phi(L_N)]}\,d\sigma^N
 =-\inf_{\mu\in\mathcal P(X)}[F_\gamma(\mu)+\Phi(\mu)].
 \label{att:full-laplace-eq}
\end{equation}
The empirical laws of the normalized attractive ensemble satisfy a
full speed-$N$ LDP with good rate
$F_\gamma-\inf F_\gamma$.
\end{theorem}

\begin{proof}
The case $\gamma=0$ is Sanov. Assume $0<\gamma<1/3$ and first take
$\Phi$ Lipschitz. Its finite-level Gibbs variational identity is
\[
 N^{-1}\log Z_{k,\gamma,\Phi}
 =\gamma\mathbb E E_k-N^{-1}D(P_{k,\gamma,\Phi}\Vert\sigma^N)
                         -\mathbb E\Phi(L_N).
\]
Every subsequence has an empirically convergent further subsequence.
Theorem~\ref{att:gibbs-energy-limit} identifies its limit as
$\int[\gamma E-H-\Phi]\,d\Pi$, at most
$\sup_\mu[\gamma E(\mu)-H(\mu)-\Phi(\mu)]$.
For the opposite bound, insert a smooth positive product law $\mu^N$
in the variational principle and use
Theorem~\ref{att:general-product-energy}.

Smooth positive references suffice for the variational infimum.
For any finite-entropy $\mu$, unitary heat averaging gives smooth
positive $\mu_\epsilon\to\mu$ with $H(\mu_\epsilon)\to H(\mu)$:
entropy contracts and is lower semicontinuous. Energy also converges,
since its unitary invariance and convexity give
$E(\mu_\epsilon)\le E(\mu)$ and its lower semicontinuity gives the
reverse lower limit. Thus the lower bound extends to all relevant
measures.

Lipschitz functions are uniformly dense in the continuous functions
on the compact metric space $\mathcal P(X)$. Both finite pressure and
the variational limit are 1-Lipschitz in the uniform norm of $\Phi$.
This proves the assertion for every continuous tilt. The full Laplace
principle on a compact probability space is equivalent to the stated
LDP.
\end{proof}

\begin{theorem}[Full attractive LDP on shrinking design tubes]
\label{att:design-laplace}
Fix any design order $t\ge1$ and $0<a<1/2$. Let
\[
 \begin{aligned}
 \mathcal D_t&=\{\mu:M_t(\mu)=I/\tbinom{t+2}{2}\},\\
 \Omega_N&=\{\|M_t(L_N)-I/\tbinom{t+2}{2}\|_{\rm HS}<N^{-a}\},
 \qquad Q_N=\sigma^N(\cdot\mid\Omega_N).
 \end{aligned}
\]
For every fixed $0\le\gamma<1/3$ and every continuous empirical tilt,
\begin{equation}
 \lim_{k\to\infty}\frac1N\log\int
 e^{N[\gamma E_k-\Phi(L_N)]}\,dQ_N
 =-\inf_{\mu\in\mathcal D_t}[H(\mu)-\gamma E(\mu)+\Phi(\mu)].
 \label{att:design-laplace-eq}
\end{equation}
The normalized empirical laws satisfy the full LDP with good rate
$F_\gamma+\iota_{\mathcal D_t}-\inf_{\mathcal D_t}F_\gamma$.
\end{theorem}

\begin{proof}
The reference probability of $\Omega_N$ tends to one by the finite
moment variance bound, since $N^{1-2a}\to\infty$.
For the upper bound, fix a closed moment tube of radius $\delta>0$.
For large $N$ it contains $\Omega_N$. The unrestricted LDP and its
bounded continuous tilt bound the upper pressure by the variational
infimum on that fixed tube. Let $\delta\downarrow0$, using compactness
and lower semicontinuity of $F_\gamma+\Phi$.

For the lower bound take a smooth positive $\mu\in\mathcal D_t$.
Under $\mu^N$, the event $\Omega_N$ has probability tending to one.
Theorem~\ref{att:general-product-energy} and its $L^1$ version show
that conditioning on this event preserves convergence of mean energy
to $E(\mu)$. Its relative entropy per particle to $Q_N$ tends to
$H(\mu)$: write out the conditioned product density, use boundedness
of $\log(d\mu/d\sigma)$, and note that both conditioning logarithms
are $o(N)$. The variational principle gives the desired lower bound.
Finally, unitary heat averaging preserves an exact projective design,
since its moment target is scalar and the heat operation is an average
of simultaneous unitary pushforwards. It approximates any finite-entropy
member of $\mathcal D_t$ with converging entropy and energy as above.
The full continuous-tilt Laplace principle yields the LDP.
\end{proof}

\begin{corollary}[Polarized manifolds and all projective dimensions]
\label{att:polarized-extension}
Let $Y$ be a compact connected complex manifold, let $(L,h)$ be a
positive Hermitian line bundle with curvature $\omega$, and let $\sigma$
be a smooth strictly positive probability volume. For all sufficiently large $k$, choose an $L^2(h^k,\sigma)$
orthonormal basis of $H^0(Y,L^k)$ and let $\Delta_k$ be its determinant,
$N_k=h^0(Y,L^k)$, and $E_k=-(kN_k)^{-1}\log\|\Delta_k\|^2$.
Write $E_\omega$ for the pluricomplex measure energy with potential
energy normalized to vanish at zero, and put
\[
 \alpha(\omega,\sigma)=\sup\left\{a>0:
 \sup_{\substack{u\in\operatorname{PSH}(Y,\omega)\\ \sup u=0}}
             \int e^{-au}\,d\sigma<\infty\right\}.
\]
For every fixed $0\le\gamma<\alpha(\omega,\sigma)$, the unrestricted
attractive determinant ensemble has the full speed-$N_k$ LDP with good
rate
\[
 H_\sigma-\gamma E_\omega-\inf(H_\sigma-\gamma E_\omega).
\]
Its Laplace principle holds for every continuous empirical-law tilt.
For $Y=\mathbb P^q$, $q\ge1$, $L=\mathcal O(q+1)$, and Fubini--Study probability,
the admissible range is $0\le\gamma<1/(q+1)$; in that range the same
shrinking-tube conclusion holds for every fixed projective design order,
with target $I/\binom{q+t}{q}$.
\end{corollary}

\begin{proof}
The proof above uses only compactness of normalized psh functions,
a uniform exponential estimate below $\alpha$, and the repulsive determinant
pressure theorem. These hold for the stated polarized manifold; the
repulsive input is \cite[Theorem 5.7]{prob_berman2017}. The good-rate
argument of Lemma~\ref{att:good-rate} uses only the same integrability
margin and the basic potential-energy duality of \cite{BBGZ2013}.
For clarity, there is no hidden homogeneity assumption in the reference
approximation. Use a smooth uniformly elliptic diffusion whose invariant
volume is $\sigma$; its heat semigroup contracts all $L^p(\sigma)$,
preserves positive lower bounds, and has a positive smooth kernel.
It therefore supplies the finite smooth reference nets. For an arbitrary
finite-entropy density its heat regularizations have converging entropy,
and their energies converge by Lemma~\ref{att:good-rate}. This replaces
the unitary averaging argument without imposing homogeneity on $Y$. The lower
Hadamard bound is $E_k\ge-o(1)$, since the smooth positive Bergman kernel
has polynomial growth in $k$. Every other step is unchanged. This proves
the unrestricted assertion; no moment constraint on a general manifold
is assumed here.

On projective space put $D=q+1$. The normalized psh threshold is
$1/D$, as can also be seen directly. Choose a maximum point of $u$ and
disintegrate Haar measure along projective lines through it. The
conditional radial density on a line is $q(1-x)^{q-1}dx$, bounded by
$q$ times line Haar density. The Riesz representation of $u/D$ on the
line, normalized at the maximum point, and Jensen's inequality give
\[
 \int_{\mathbb P^q}e^{-au}\,d\sigma\le\frac{q}{1-Da},
 \qquad 0\le a<1/D.
\]
Indeed a unit linear form has squared modulus uniform on $[0,1]$ under
line Haar measure. Conversely $u=D\log|z_0|^2$ has supremum zero and
fails exponential integrability at $a=1/D$, proving sharpness of this
one-coordinate threshold. Finally unitary heat averaging on
$\mathbb P^q$ preserves each exact projective design. The proof of
Theorem~\ref{att:design-laplace} applies verbatim in every dimension.
\end{proof}

\begin{corollary}[Zero pressure and the unique Haar equilibrium]
\label{att:haar-minimizer}
On $\mathbb P^q$ with anticanonical polarization, let
$0\le\gamma<1/(q+1)$. Then
\[
 \lim_k N_k^{-1}\log Z_{k,\gamma}=0,
 \qquad H(\mu)-\gamma E(\mu)\ge0,
\]
and equality in the second assertion holds exactly for $\mu=\sigma$.
The unconditioned attractive rate is therefore $F_\gamma$ itself.
For every fixed-order shrinking design tube the rate is
$F_\gamma+\iota_{\mathcal D_t}$, and its unnormalized pressure is also
zero.
\end{corollary}

\begin{proof}
Corollary~\ref{ac:tubes} gives empirical convergence to $\sigma$.
The entropy and energy limit in Theorem~\ref{att:gibbs-energy-limit},
with the same proof on projective space of every dimension as explained
in Corollary~\ref{att:polarized-extension}, then gives
\[
 N^{-1}D(P_{k,\gamma}\Vert\sigma^N)\longrightarrow0,
 \qquad \mathbb E_{P_{k,\gamma}}E_k\longrightarrow E(\sigma)=0.
\]
The finite Gibbs identity yields zero pressure. The full Laplace principle
therefore identifies $\inf F_\gamma=0$, attained by $\sigma$.
For $\gamma=0$ these assertions are also immediate from entropy.

To prove uniqueness, let $\mu\ne\sigma$. Choose a real smooth $f$
with $\mu(f)\ne\sigma(f)$, and a weak neighborhood $U$ of $\mu$
on which $|\nu(f)-\sigma(f)|>s$ for some $s>0$. The exponential
bound of Theorem~\ref{ac:concentration} gives
$P_{k,\gamma}(L_N\in U)\le2e^{-cN}$ with $c>0$.
The LDP lower bound on this open set gives
$-\inf_U F_\gamma\le-c$. In particular $F_\gamma(\mu)\ge c>0$.
The use of the exponential estimate is essential: convergence in
probability alone would not imply strict positivity of the rate away
from $\sigma$. Finally $\sigma$ is an exact design of every order,
so restriction to a design fibre preserves the unique zero and the
normalization in Theorem~\ref{att:design-laplace}.
\end{proof}

\paragraph{Range and the next problem.}
The proven interval is a fixed positive attractive interval,
$0\le\gamma<1/3$, for the natural complex determinant and for every
fixed-order shrinking design constraint. It includes all continuous
empirical-law tilts. The scalar section threshold is used twice:
for integrability and compactness of attractive conditional densities.
A result near $\gamma=1$ requires a replacement for that conditional
compactness mechanism; first-moment macroscopic coercivity alone does
not provide it. The toric assignment formula separately makes the product energy
computable by finite optimal assignments, while the polarized extension
shows that the entropy-and-conditional-compactness mechanism is not
specific to a homogeneous projective space.
